\documentclass[a4paper]{article}

\usepackage{graphicx}
\usepackage{amsmath,amssymb}
\usepackage{minted}
\usepackage{multirow}
\usepackage{float}
\usepackage{todonotes}
\usepackage{forest}
\usepackage{xstring}
\usepackage{xcolor}
\usepackage[most]{tcolorbox}
\usepackage{xparse}
\usepackage{natbib}

\usetikzlibrary{graphs,graphdrawing}
\usegdlibrary{layered}

% for external graphviz
\input{/home/donkarlo/Dropbox/repo/nd_ai_project/src/nd_ai/knoweldge_representation/language/formal/computer/programming/tex/latex/code_clips/external_graphviz.tex}

% for tags
\input{/home/donkarlo/Dropbox/repo/nd_ai_project/src/nd_ai/knoweldge_representation/language/formal/computer/programming/tex/latex/parts/control_sequence/kinds/semtag/semtag.tex}

% for links
\input{/home/donkarlo/Dropbox/repo/nd_ai_project/src/nd_ai/knoweldge_representation/language/formal/computer/programming/tex/latex/code_clips/seehere.tex}

\title{Population spike}
\date{}

\begin{document}
    \maketitle
    \clickurlfooter{https://en.wikipedia.org/wiki/Population_spike}
    
    \cite{mcnaughton-1978-synaptic-enhancement-in-fascia-dentata-cooperativity-among-coactive-afferents} studies synaptic enhancement in the fascia dentata of the hippocampus and emphasizes cooperativity among simultaneously active afferents.
    It is relevant to population spikes because that phenomenon is commonly observed in hippocampal field recordings where many nearby neurons fire in a coordinated way.
    Conceptually, the paper helps connect collective extracellular signals to underlying synaptic input and coordinated neuronal discharge.
    So, for your purpose, it is a useful early reference if you want a source tied to hippocampal population-level spiking rather than to a single neuron’s action potential
    \bibliographystyle{plainnat}
    \bibliography{/home/donkarlo/Dropbox/repo/nd_shared_research/bibliography/bibliography.bib}
\end{document}